% Generated draft: requires mathematical and source review.
\documentclass[11pt]{article}
\usepackage[margin=25mm]{geometry}
\usepackage{fontspec}
\setmainfont{Latin Modern Roman}
\usepackage{amsmath,amssymb,mathtools,unicode-math}
\setmathfont{Latin Modern Math}
\usepackage{enumitem,needspace,xurl,hyperref}
\setlength{\parindent}{0pt}
\setlength{\parskip}{5pt}
\setlength{\emergencystretch}{3em}
\newcommand{\N}{\mathbb{N}}
\newcommand{\Z}{\mathbb{Z}}
\newcommand{\Q}{\mathbb{Q}}
\newcommand{\R}{\mathbb{R}}
\newcommand{\C}{\mathbb{C}}
\newcommand{\F}{\mathbb{F}}
\newcommand{\A}{\mathbb{A}}
\newcommand{\PP}{\mathbb{P}}
\newcommand{\E}{\mathbb{E}}
\newcommand{\eps}{\varepsilon}
\newcommand{\dd}{\,\mathrm{d}}
\begin{document}
\section*{UnsolvedMath ID: 800482}
\textbf{Source problem code:} catalogue entry
\subsection{Definitions and mathematical statement}
A domain is a nonempty connected open set. A domain \(\Omega\) in a Euclidean space is convex if \((1-t)x+ty\in\Omega\) for all \(x,y\in\Omega\) and \(t\in[0,1]\). Its boundary is \(C^1\)-smooth if it is locally the graph of a continuously differentiable function.

Write \(T_\Omega\) for the symplectic billiard dynamics associated with \(\Omega\), using the ambient symplectic space, phase space, and billiard convention intended by the source. These are not defined in the supplied record. The source asks:
\[
\text{Does there exist a convex domain }\Omega\text{ with }C^1\text{-smooth boundary such that }T_\Omega\text{ is chaotic?}
\]
The term \emph{chaotic} is not defined in the supplied record. Thus this is a formulation template requiring a specified dynamical property before it becomes a precise existence statement. The ambient dimension, any boundedness requirement, and the definition and well-definedness of the intended dynamics at the stated regularity also require clarification; no additional hypotheses are imposed here.

The source identifies the stadium as a natural candidate and reports numerical evidence suggesting chaotic behavior. It supplies neither a precise chaos criterion nor details of that numerical evidence. This suggestion is not asserted as a mathematical result.
\subsection{Short English statement}
The question asks whether a convex domain with a \(C^1\)-smooth boundary can have chaotic symplectic billiard dynamics. The stadium is suggested as a candidate on the basis of reported numerical evidence. A precise statement still requires the intended definition of chaos and the ambient and dynamical conventions.
\subsection{Sources}
\begin{itemize}
\item Local UnsolvedMath registry, \texttt{data/unsolvedmath/problems.json} and the supplied additions record: permanent ID \(800482\), source problem code \texttt{catalogue entry}; title: \emph{Billiards and Quantitative Symplectic Geometry --- Serge Tabachnikov --- Problem 4}. The registry attributes the collection to Albach, Barraud et al., with proposed year 2026. Its supplied status is \emph{open}; this label is not an independently verified mathematical or literature assessment.
\item Albach, Barraud et al., \emph{Open Problems in Billiards and Quantitative Symplectic Geometry} (2026), source item \emph{Serge Tabachnikov --- Problem 4}; contributor: Serge Tabachnikov. Supplied source URL: \url{https://arxiv.org/abs/2602.12896}. The record reports access on 2026-10-16 and extraction by \texttt{arxiv\_latex}. These are supplied provenance details; the source was not independently accessed or verified for this draft.
\end{itemize}
\end{document}
